Mira las tablas \ref{tab:miTablaSinBordes}, \ref{tab:miTablaBordesVerticales} y \ref{tab:miTablaBordesHorizontales}

\begin{table}[h]
    \centering
    \begin{tabular}{cccc}   % l, c, r -> Alineación de cada columna
        1 & 2 & 3 & 4 \\
        a & b & c & d\\
        g & h & i & j\\
    \end{tabular}
    \caption{Tabla sin bordes}
    \label{tab:miTablaSinBordes}
\end{table}

\begin{table}[h]
    \centering
    \begin{tabular}{|c|c|c|c|}  % Aquí indicamos lso bordes verticales
        1 & 2 & 3 & 4 \\
        a & b & c & d\\
        g & h & i & j\\
    \end{tabular}
    \caption{Tabla con bordes verticales}
    \label{tab:miTablaBordesVerticales}
\end{table}

\begin{table}[h]
    \centering
    \begin{tabular}{|c||c|c|c|}  % Aquí indicamos lso bordes verticales, también duplicados si queremos
        \hline
        1 & 2 & 3 & 4 \\
        \hline      % Ponemos doble borde
        \hline
        a & b & c & d\\
        \hline
        g & h & i & j\\
        \hline
    \end{tabular}
    \caption{Tabla con bordes horizontales}
    \label{tab:miTablaBordesHorizontales}
\end{table}
